%%% Dokumentation der Nutzung von KI-Tools (Kurzfassung aus Textbausteinen)
%%% Zutreffende Textbausteine verwenden, Unzutreffendes bitte löschen.
%%% Weitere Arten der Nutzung, die auf dieser Liste nicht erscheinen, sind
%%% unbedingt vorab mit Ihrem/Ihrer Betreuer/in zu klären.

\begin{aiusage}
	%\aiusednone                          % keinerlei Hilfsmittel der generativen KI benutzt
	%\aiusedfor{ideas}                    % um erste Ideen für Forschungsfragen und Theorien zu sammeln
	\aiusedfor{literature-search}         % um Literatur zu finden
	%\aiusedfor{literature-understanding} % um Literatur besser nachvollziehen zu können
	\aiusedfor{proofreading}              % um Korrektur zu lesen und den sprachlichen Stil zu verbessern
	%\aiusedforother{um ... zu ...}       % weitere Anwendungen
\end{aiusage}

%%% Alternativ: ausführliche Dokumentation (Langfassung), falls vom Studiengang verlangt.
%%% Aktivitäten für \aiactivity{<Aktivität>}{<KI-Tools>}{<Beschreibung der Nutzung>}:
%%% ideas, literature, methodology, coding, data, interpretation, structure,
%%% generation, translation, editing, references

%\begin{aiusagedetails}
%	\aiconfirm{accuracy}         % Ergebnisse der KI-Tools überprüft und ggf. korrigiert
%	\aiconfirm{responsibility}   % Verantwortung liegt bei den Autor*innen
%	\aiactivity{literature}{Perplexity}{Suche nach Literatur zu ... (Kapitel~2)}
%	\aiactivity{editing}{DeepL Write}{Korrekturlesen der Kapitel~1--5}
%	%\aifurtheruse{...}
%\end{aiusagedetails}
